\begin{table}[htbp]
\centering
\caption{Dez maiores hiatos absolutos de renda domiciliar per capita segundo sexo da pessoa responsável, Sergipe, 2022}
\label{tab:item2-4-maiores-hiatos-rdpc}
\small
\begin{tabular}{rlrrr}
\toprule
Posição & Município & Resp. mulher (R\$) & Resp. homem (R\$) & Hiato (R\$) \\
\midrule
1 & Aracaju & 2.019,85 & 2.809,23 & 789,37 \\
2 & Barra dos Coqueiros & 1.486,98 & 2.138,42 & 651,44 \\
3 & São Cristóvão & 1.123,64 & 1.505,71 & 382,07 \\
4 & Rosário do Catete & 719,48 & 1.074,16 & 354,68 \\
5 & Santo Amaro das Brotas & 788,98 & 1.111,90 & 322,92 \\
6 & Carmópolis & 769,38 & 1.075,72 & 306,34 \\
7 & Itabaiana & 1.070,28 & 1.367,85 & 297,57 \\
8 & Nossa Senhora do Socorro & 806,32 & 1.095,68 & 289,36 \\
9 & Carira & 721,61 & 1.002,37 & 280,76 \\
10 & Itaporanga d'Ajuda & 663,94 & 941,88 & 277,93 \\
\bottomrule
\end{tabular}
\fonte{Elaboração própria com base nos microdados do Censo Demográfico 2022.}
\nota{Hiato calculado como renda per capita média de domicílios com responsável homem menos a de domicílios com responsável mulher.}
\end{table}
